\section{Related Work}

\paragraph{Pre-trained text representations}
Pre-trained textual representations are widely used to improve performance on NLP tasks.
These representations are trained on large corpora (usually unsupervised), and fed as features to downstream models.
In deep networks, these features may also be fine-tuned on the downstream task.
Brown clusters, trained on distributional information, are a classic example of pre-trained representations~\citep{brown1992}.
\citet{turian2010} show that pre-trained embeddings of words outperform those trained from scratch.
Since deep-learning became popular, word embeddings have been widely used, and many training strategies have arisen~\citep{mikolov2013,pennington2014,bojanowski2017enriching}.
Embeddings of longer texts, sentences and paragraphs, have also been developed~\citep{le2014,kiros2015,conneau2017,cer2019}.

To encode context in these representations, features are extracted from internal representations of sequence models,
such as MT systems~\citep{mccann2017}, and BiLSTM language models, as used in ELMo~\citep{peters2018}.
As with adapters, ELMo exploits the layers other than the top layer of a pre-trained network.
However, this strategy only \emph{reads} from the inner layers.
In contrast, adapters \emph{write} to the inner layers, re-configuring the processing of features through the entire network.

\paragraph{Fine-tuning}
Fine-tuning an entire pre-trained model has become a popular alternative to features~\citep{dai2015,howard2018universal,radford2018improving}
In NLP, the upstream model is usually a neural language model~\citep{bengio2003}.
Recent state-of-the-art results on question answering~\citep{rajpurkar2016} and text classification~\citep{wang2018glue} have been attained by fine-tuning a Transformer network~\citep{vaswani2017} with a Masked Language Model loss~\citep{devlin2018bert}.
Performance aside, an advantage of fine-tuning is that it does not require task-specific model design, unlike representation-based transfer.
However, vanilla fine-tuning does require a new set of network weights for every new task.

\paragraph{Multi-task Learning}
Multi-task learning (MTL) involves training on tasks simultaneously.
Early work shows that sharing network parameters across tasks exploits task regularities, yielding improved performance~\citep{caruana1997}.
The authors share weights in lower layers of a network, and use specialized higher layers.
Many NLP systems have exploited MTL.
Some examples include: text processing systems (part of speech, chunking, named entity recognition, etc.)~\citep{collobert2008}, multilingual models~\citep{huang2013cross}, semantic parsing~\citep{peng2017}, machine translation~\citep{johnson2017}, and question answering~\citep{choi2017}.
MTL yields a single model to solve all problems.
However, unlike our adapters, MTL requires simultaneous access to the tasks during training.

\paragraph{Continual Learning}
As an alternative to simultaneous training, continual, or lifelong, learning aims to learn from a sequence of tasks~\citep{thrun1998}.
However, when re-trained, deep networks tend to forget how to perform previous tasks; a challenge termed catastrophic forgetting~\citep{mccloskey1989catastrophic,french1999catastrophic}.
Techniques have been proposed to mitigate forgetting~\citep{kirkpatrick2017overcoming,zenke2017continual}, however, unlike for adapters, the memory is imperfect.
Progressive Networks avoid forgetting by instantiating a new network ``column'' for each task~\citep{rusu2016progressive}.
However, the number of parameters grows linearly with the number of tasks,
since adapters are very small, our models scale much more favorably.

\paragraph{Transfer Learning in Vision}
Fine-tuning models pre-trained on ImageNet~\citep{deng2009} is ubiquitous when building image recognition models~\citep{yosinski2014,huh2016makes}.
This technique attains state-of-the-art performance on many vision tasks, including classification~\citep{kornblith2018better}, fine-grained classifcation~\citep{hermans2017}, segmentation~\citep{long2015}, and detection~\citep{girshick2014}.
In vision, convolutional adapter modules have been studied~\citep{rebuffi2017,rebuffi2018,rosenfeld2018incremental}.
These works perform incremental learning in multiple domains by adding small convolutional layers to a ResNet~\citep{he2016} or VGG net~\citep{simonyan2014very}.
Adapter size is limited using $1\times 1$ convolutions, whilst the original networks typically use $3\times 3$.
This yields $11\%$ increase in overall model size per task.
Since the kernel size cannot be further reduced other weight compression techniques must be used to attain further savings.
Our bottleneck adapters can be much smaller, and still perform well.

Concurrent work explores similar ideas for BERT~\citep{stickland2019bert}.
The authors introduce Projected Attention Layers (PALs), small layers with a similar role to our adapters.
The main differences are i) \citet{stickland2019bert} use a different architecture,
and ii) they perform multitask training, jointly fine-tuning BERT on all GLUE tasks.
\citet{semnani2019} perform an emprical comparison of our bottleneck Adpaters and PALs on SQuAD v2.0~\citep{rajpurkar2018}.
